%───────────────────────────────────────────────────────────────────────────────
\subsection{v2 Training Improvements}
\label{sec:v2}

Based on the v1 training diagnostics (28{,}980 optimiser steps, 690 epochs),
we identified three failure modes and designed targeted improvements.  All
changes are hyperparameter tweaks or minor code modifications; no architectural
changes are introduced.

\textbf{1.\ Boundary-aware time sampling} (Section~\ref{sec:time_sampling}).
The logit-normal distribution assigns $< 0.04\%$ of training samples to the
1-NFE region $(t \approx 1, r \approx 0)$.  We inject $\alpha = 10\%$ of each
batch from $\mathcal{U}([0.95, 1] \times [0.001, 0.05])$ to provide direct
training signal at the 1-NFE operating point.

\textbf{2.\ Extended training with improved checkpointing.}
Max epochs increased from 1{,}500 to 3{,}000 (${\sim}$126K steps,
$4.3\times$ v1 budget).  Early-stopping patience raised from 10 to 30 FID
evaluations, with FID computed every validation epoch (every 10 training epochs)
instead of every 3rd.  The best-FID checkpoint callback retains the top-3
models (previously top-1) for rollback and post-hoc analysis.

\textbf{3.\ Learning rate warmup.}
Linear warmup over 1{,}000 steps (${\sim}$24 epochs), addressing the 90.5\%
gradient clipping fraction observed at epoch~0.  Justified by the large
effective batch size ($B_\text{eff} = 132$) per Goyal \etal~\cite{goyal2017accurate}.

\textbf{4.\ Inference-time norm correction.}
Post-training calibration of a scalar $\gamma$ (Section~\ref{sec:1nfe})
to compensate for compound velocity norm overshoot ($\|V_\theta\|/\|v_c\|
= 1.16$ at epoch~689).

\textbf{5.\ Per-time-bin cosine diagnostics.}
Cosine similarity $\cos(V, v_c)$ and $\cos(\tilde{v}, v_c)$ logged per
time bin (near-data, mid, near-noise) for finer-grained convergence monitoring.
